\chapter[Introdução]{Introdução}
\label{cap:introducao}

\section[Contextualização]{Contextualização}
\label{sec:contextualizacao}

Um importante aspecto para o desenvolvimento de um país é a educação. Através dela, muda-se a vida tanto de pessoas em sua individualidade quanto no coletivo. Qualificação de mão de obra, novas descobertas tecnológicas, entre outros pontos positivos tornam importante o investimento e desenvolvimento de uma boa educação para os cidadãos de uma nação, e o Brasil não é exceção.

O Índice de Desenvolvimento da Educação Básica (Ideb) é um indicador da qualidade da educação brasileira criado em 2007, visando apurar a qualidade dos ensinos fundamental (separado em anos iniciais do 1° ao 5° ano e em anos finais do 6° ao 9° ano) e médio. Seu resultado leva em conta duas medições: os dados sobre aprovação escolar, obtidos no Censo Escolar, e as médias de desempenho no Sistema de Avaliação Básica (Saeb). O índice varia de 0 a 10, sendo que um valor 6 é equivalente ao sistema educacional de um país desenvolvido \cite{inep2024indicador}.

Entretanto, nota-se que chegar a tal valor ainda é um desafio. Em 2023 o Brasil alcançou nota 6 nos anos iniciais ensino fundamental (1° ao 5° ano), 5 nos anos finais fundamental (6° ao 9° ano) e 4,3 no ensino médio, aumentando apenas 0,1 porcento se comparado a anos anteriores. Em todos os casos, o ensino público sempre obtém pior desempenho \cite{ferreira2024educacao}. 

Existem diferentes metodologias de aprendizagem, sendo algumas delas de aprendizagem colaborativa, onde estudantes se ajudam em seu processo de aprendizagem. Uma delas é a monitoria estudantil, em que um aluno já aprovado em uma disciplina auxilia outros estudantes na aprendizagem dela. Nos cursos superiores, a monitoria tem sido utilizada com muita frequência como estratégia de apoio ao ensino, especialmente para atender estudantes com dificuldades de aprendizagem. Ao atuar como monitor(a), o estudante auxilia o(a) professor(a) no processo de ensino-aprendizagem, aprofunda seus conhecimentos na determinada disciplina e propicia ao graduando desenvolver o interesse pela carreira docente, pois convive com a prática diária do ensino \cite{gonccalves2021importancia}.

Uma metodologia similar é a Trezentos, criada por Ricardo Fragelli. Uma avaliação individual é aplicada na turma de uma disciplina e, após os resultados, estudantes são agrupados em grupos de cinco a seis membros, de maneira que os com alto rendimento se unem aos com baixo rendimento, e os auxiliam na realização de metas estabelecidas pelo professor para que os estudantes com baixa nota do grupo tenham a chance de refazer a avaliação \cite{trezentos}. Em 2015, o índice de aprovação na disciplina Cálculo 1 da Universidade de Brasília (UnB) no campus Gama com a aplicação do Trezentos  subiu de 50\% para 95\% \cite{trezentosResultados}. Essas duas metodologias provam que, de fato, a aprendizagem colaborativa pode ser muito eficaz.

Percebe-se que, de forma geral, há uma crise motivacional, principalmente no que tange ao cenário educacional. Grande parte das instituições de ensino, independente de nacionalidade e de níveis de educação, encontra dificuldades para engajar seus alunos utilizando os recursos educacionais tradicionais. Uma metodologia recente que ajuda a mitigar o problema é a gamificação, que consiste na utilização de elementos de jogos fora do contexto deles, para promover a aprendizagem, motivando os aprendizes a alguma ação e auxiliando na solução de problemas e na interação com outros estudantes  \cite{tolomei2017gamificaccao}.

Elementos da gamificação incluem sistemas de pontuação, níveis, ranking, medalhas, desafios, entre outros recursos mais comuns aos jogos. Trabalhando com pontuações e níveis de experiência, o usuário é instigado a buscar atividades a fim de cumprir metas e atingir objetivos. Esses fatores dialogam entre si, aumentando o senso de socialização e colaboração, sem contar o aumento de \textit{feedback} contínuo, proporcionando a noção do progresso em uma atividade realizada durante o processo de aprendizagem \cite{tolomei2017gamificaccao}.

Diante dos fatores citados, um software que forneça a oportunidade de intermediar o contato entre pessoas que dominam uma área do saber e estudantes interessados em aprender e que utilize elementos de gamificação para manter o engajamento de ambas as partes tem potencial de auxiliar na aprendizagem de estudantes com dificuldade de aprendizagem.


\section[Questão de Pesquisa]{Questão de Pesquisa}
\label{sec:questao}

Com o objetivo de nortear o desenvolvimento do trabalho, a seguinte questão de
pesquisa foi definida:

\textit{Como uma plataforma pode auxiliar estudantes com dificuldades em uma área de conhecimento e pessoas versadas nesta área a se conectarem e agendarem encontros ao mesmo tempo que os mantém engajados a usá-la por tempo prolongado, aumentando a assimilação e fixação de conhecimento?}

\section[Justificativa]{Justificativa}
\label{sec:justificativa}
% Por que? Para melhorar a capacidade de aprendizagem

A situação da educação é um problema de longa data no Brasil. Além dos resultados nacionais do Ideb mencionados anteriormente, o desempenho do país se encontra entre as piores posições em avaliações internacionais, como por exemplo o \textit{Programme for International Student Assesment} (PISA) na educação básica e \textit{QS-World University Ranking} (QS-WUR) na educação superior.

O Brasil obteve as seguintes posições no ano de 2022: entre 62ª e 69ª posição no PISA em matemática, entre 44ª e 57ª posição em leitura e entre 53ª e 64ª posição em ciências \cite{pisa2024}.
Enquanto na lista de 92ª posição  no QS-WUR, em que apenas a Universidade de São Paulo (USP) aparece entre as universidades do pais no top 100 do ranking, as demais universidades do Brasil ficam em posições maiores de 200 \cite{qswur2025}.

Não é incomum que cada vez mais alunos se encontrem desinteressados e dispersos pelo conteúdo em sala de aula, principalmente com as metodologias de aprendizagem, as quais são majoritariamente passivas \cite{fragelli2017game}.

Um fenômeno emergente é a gamificação, que deriva da popularidade dos \textit{games}, e da capacidade intrínseca de motivar ação, resolver problema e melhorar aprendizagem nas diversas áreas do conhecimento humano \cite{fardo2013gamificacao}.
No campo educacional, a gamificação tem sido utilizada cada vez mais, superando
os métodos tradicionais de ensino e aprendizagem pelos games \cite{souza2025metodologias}.

A gamificação é a infusão intencional de elementos e mecânicas de jogos em situações que não são de jogos. O objetivo principal é engajar os participantes ativamente, motivá-los e ajudá-los a alcançar objetivos específicos. Essa abordagem aproveita nossas tendências naturais, como o desejo de competição, o senso de realização, o trabalho em equipe e o altruísmo, para impulsionar a participação, o desempenho e a satisfação geral \cite{magtani2024design}.

Há vários artigos com relatos de experiências de professores ao redor do mundo que utilizaram esses métodos em diversas áreas do conhecimento, da educação infantil ao ensino superior, obtendo resultados positivos \cite{fragelli2017game}.

Beat the GMAT (BTG) é a maior rede social do mundo para os candidatos de MBA (\textit{Masters in Business Administration}), atendendo a mais de 2 milhões de pessoas a cada ano. A BTG habilita as pessoas a aprender, compartilhar, ensinar e apoiar uns aos outros ao longo do processo de admissão de MBAs de grandes universidades. Essa empresa adotou a gamificação para aumentar o engajamento e retenção de seus usuários \cite{tolomei2017gamificaccao}.

A BTG foi configurada para, com base em determinados comportamentos, o usuário ser recompensado e ter em seu perfil essas recompensas. Um exemplo é a medalha ``Campeão da Gramática'' para usuários que escrevem 100 postagens no fórum de correção de gramática \cite{tolomei2017gamificaccao}.

Com essas e outras ações gamificadas, a BTG aumentou o tempo gasto por seus usuários em seu site e em sua comunidade em 370\%, significando um engajamento maior dos usuários nas atividades propostas. Além disso, a proposta gamificada trouxe aumento de 195\% de visitação em suas páginas, proporcionando à empresa uma possibilidade maior de conversão em novos usuários \cite{tolomei2017gamificaccao}.

Diante desses fatos, o tema deste trabalho se mostrou muito interessante, pois propõe unir duas abordagens educacionais em uma única plataforma, sendo potencialmente útil para a aprendizagem de alunos com problemas educacionais.

% O que? Plataforma de tutoria online

% Desta forma, este trabalho visa propor o desenvolvimento de uma plataforma de tutoria online, com elementos de gamificação, com o intuito mitigar o problema com a educação do Brasil. Incentivando tutores e alunos a se manterem engajados em ensinar e aprender.

% Quem? Tutores e alunos
% Qual?
% Aspectos positivos
% Aspectos impactos

\section[Objetivos]{Objetivos}
\label{sec:objetivos}

\subsection[Objetivo geral]{Objetivo Geral}

Este trabalho tem como objetivo geral propor a implementação de uma aplicação web que intermedia o contato entre tutores de uma área de conhecimento e pessoas interessadas nesta tutoria.

\subsection[Objetivos específicos]{Objetivos Específicos}

\begin{itemize}
    \item Estruturar um ambiente virtual que otimize a busca e conexão entre aprendizes e tutores, utilizando filtros parametrizados para facilitar a localização de tutores na área de conhecimento desejada;
    
    \item Aplicar técnicas de sistemas de recomendação para personalizar a sugestão de tutores, visando mitigar a sobrecarga de escolha e aumentar a relevância dos resultados para o aprendiz;
    
    \item Utilizar estratégias de gamificação para estimular a motivação e sustentar o engajamento de longo prazo de tutores e aprendizes na plataforma;

    \item Sistematizar o processo de agendamento de sessões de tutorias para assegurar a conciliação de horários e a organização das sessões;
    
    \item Viabilizar mecanismos de interação e \textit{feedback} direto entre os tutores e os aprendizes para consolidar a confiança e a construção de reputação dentro da comunidade da aplicação.
    
\end{itemize}

\section[Metodologias]{Metodologias}
\label{sec:metodologia}

As metodologias deste trabalho podem ser divididas em metodologia de pesquisa e metodologia de desenvolvimento.

Uma metodologia de pesquisa é caracterizada por sua abordagem, natureza, objetivo e procedimentos \cite{gerhardt2009metodos}. A metodologia de pesquisa deste trabalho tem abordagem qualitativa, natureza de pesquisa aplicada, objetivo de pesquisa exploratória e pesquisa bibliográfica como procedimento.

A metodologia de desenvolvimento de software utilizada consiste em uma adaptação da metodologia \textit{eXtreme Programming} (XP), com adição da forma de acompanhamento do \textit{Kanban}, levando em consideração o escopo do projeto, e também considerando apenas os dois alunos como desenvolvedores. Essas metodologias serão melhores explicadas no Capítulo \ref{cap:metodologia}.

\section[Organização do Trabalho]{Organização do Trabalho}
\label{sec:organizacao}
Este trabalho foi dividido em cinco capítulos, sendo esses resumidos da seguinte maneira:

\begin{itemize}
    \item Capítulo \ref{cap:introducao}: este capítulo introduz o tema do trabalho, compreendendo assim à \nameref{sec:contextualizacao}, \nameref{sec:questao}, \nameref{sec:justificativa}, \nameref{sec:objetivos} e \nameref{sec:metodologia}.

    \item Capítulo \ref{cap:referencial}: este capítulo sintetiza os principais conhecimentos teóricos que fornecem a melhor compreensão do contexto no qual o trabalho está envolvido, sendo este dividido em quatro seções: \nameref{sec:tutoria}, \hyperref[sec:aw]{Aplicação Web}, \hyperref[sec:recomendação]{Sistemas de Recomendação} e \hyperref[sec:gamificacao]{Gamificação}.

    \item Capítulo \ref{cap:metodologia}: este capítulo apresenta a proposta do trabalho, esclarecendo os principais aspectos quanto à metodologia utilizada, tecnologias escolhidas, requisitos definidos, arquitetura de software e desenvolvimento do trabalho proposto.
    
    \item Capítulo 4: este capítulo apresenta a execução da proposta, bem como as mudanças necessárias para o sucesso da implementação.

    \item Capítulo 5: este capítulo finaliza o projeto proposto, apresentando os resultados alcançados e trabalhos futuros relacionados ao projeto.
\end{itemize}